\documentclass{article}
\usepackage[a4paper,margin=2.8cm]{geometry}

\usepackage{lmodern}
\usepackage[french]{babel}
\usepackage[T1]{fontenc}
\usepackage{mathtools, amssymb,amsthm}


\usepackage{url,mdframed}
\usepackage{xcolor}

\usepackage{hyperref}
\hypersetup{
    colorlinks,
    linkcolor={red!50!black},
    citecolor={blue!50!black},
    urlcolor={blue!80!black}
}

\def\separateur{\begin{center}
$\star$\\
$\star\quad\star$
\end{center}}

\begin{document}
\author{Damien Mégy}
\title{Bibliographie commentée\\M1 de mathématiques \\ \&\\ Deuxième année de Magistère de mathématiques}

\maketitle
\tableofcontents

% il manque les miniatures, proof from the book

\section*{Comment utiliser cette bibliographie}

\newpage
\section{Bibliographie pour le M1}

\subsection{Algèbre}

\subsection{Géométrie et topologie}

\subsection{Probabilités et statistiques}

\subsection{Analyse et équations aux dérivées partielles}





\end{document}

A ajouter:

The Art and Craft of Problem Solving de Paul Zeitz

Mathematics: A Very Short Introduction de Timothy Gowers

Mathematics: Its Content, Methods and Meaning de Aleksandrov, Kolmogorov, Lavrent’iev

Letters to a Young Mathematician de Ian Stewart

AUtres livres de Stewart ?



\section{Bibliographie spécifique pour les cours supplémentaires de Magistère}

\subsection{Analyse harmonique et applications à la théorie des nombres}
\begin{itemize}
\item  \emph{}
\begin{quote}

\end{quote}
\item  \emph{}
\begin{quote}

\end{quote}
\item  \emph{}
\begin{quote}

\end{quote}
\end{itemize}

\subsection{Mécanique quantique}
\begin{itemize}
\item L.D. Landau, E. M. Lifshitz, \emph{Mécanique quantique, théorie non relativiste}
\begin{quote}
Troisième tome du Landau-Lifshitz, mythique ouvrage soviétique.
\end{quote}
\item R. Feynman, R. Leighton, M. Sands, \emph{Mécanique quantique}
\begin{quote}
Le tome 5 tome du \og Cours de physique de Feynman\fg. (Feynman, professeur d'université à 24 ans, prix Nobel de physique et grand vulgarisateur.) 
\end{quote}
\item C. Cohen-Tannoudji, \emph{Mécanique quantique}
\begin{quote}
Autre bible de mécanique quantique, cette fois par un prix Nobel made in France.
\end{quote}
\end{itemize}

\subsection{Chaos et systèmes dynamiques}
\begin{itemize}
\item  \emph{}
\begin{quote}

\end{quote}
\item  \emph{}
\begin{quote}

\end{quote}
\item  \emph{}
\begin{quote}

\end{quote}
\end{itemize}

\subsection{Logique et modèles de calcul}
\begin{itemize}
\item  \emph{}
\begin{quote}

\end{quote}
\item  \emph{}
\begin{quote}

\end{quote}
\item  \emph{}
\begin{quote}

\end{quote}
\end{itemize}

\subsection{Analysis Situs : introduction à la topologie algébrique}
\begin{itemize}
\item  \emph{}
\begin{quote}

\end{quote}
\item  \emph{}
\begin{quote}

\end{quote}
\item  \emph{}
\begin{quote}

\end{quote}
\end{itemize}

